% RESULTADO_RECUPERADO. II REV10,
% 05b_coordenadas_angulares.tex:1--125.
% Se conserva la publicación angular, sus cartas y pruebas; el funcional
% ulterior de Barbero no es una dependencia de la forma elíptica.
\section{Représentation angulaire et orientation des canaux}
\label{iv:angular:section}

La section \(\hbar_{\mathrm{ret}}\) appartient à la droite
d’action \(\mathcal L_S\) ; \(\eta_{\mathrm{ret}}\) est son
coefficient sans dimension dans la base déclarée dans
\eqref{iv:action:hbar}. La représentation angulaire agit ensuite
sur \((\alpha,\eta_{\mathrm{ret}})\). En écrivant
\([t]_{360}\) pour la classe de \(t\) dans
\(\mathbb R/(360\mathbb Z)\), on définit
\begin{equation}
 \mathcal P_{\mathrm{ang}}(\alpha,\eta_{\mathrm{ret}})
 =\bigl([1000\alpha]_{360},
        [2(\eta_{\mathrm{ret}}+\alpha)]_{360}\bigr).
 \label{iv:angular:publicacion}
\end{equation}
La complétion nonadique détermine le facteur \(1000\) ; la
clôture bidirectionnelle détermine le facteur deux. Le système de coordonnées circulaire
à \(360\) unités conserve les partitions
\(12\cdot30=4\cdot90=3\cdot120\). Sur la branche positive, les
représentants relevés sont
\begin{equation}
 A_{\deg}=1000\alpha,\qquad
 C^*_{\deg}=2(\eta_{\mathrm{ret}}+\alpha).
 \label{iv:angular:representantes}
\end{equation}
L’indice distingue ces coordonnées angulaires des blocs
du registre dodécaphasique.

\begin{proposition}[Décalage de la représentation en base mille]
\label{iv:angular:desplazamiento}
Soit \(\alpha=\sum_{n\geq1}b_n1000^{-n}\), avec
\(b_n\in\{0,\ldots,999\}\), le développement compatible déjà
engendré. Sa coordonnée relevée satisfait
\[
 A_{\deg}=b_1+\sum_{n\geq1}b_{n+1}1000^{-n}.
\]
La représentation conserve les blocs à partir du deuxième et
place le bloc initial dans la partie entière.
\end{proposition}
\begin{proof}
La série converge absolument. Multiplier chaque terme par
\(1000\) et séparer le terme initial fournit l’identité.
L’opération agit sur le développement construit, avec ses
blocs et sa mémoire ; la représentation circulaire ultérieure prend
sa classe modulo \(360\).
\end{proof}

\subsection{Systèmes de coordonnées circulaires et covariance de l’action}
\label{iv:angular:cartas}

\begin{definition}[Coordinatisations de la coordonnée angulaire]
\label{iv:angular:unidades}
Le système de coordonnées en degrés exprime un tour complet au moyen de \(360\)
unités ; le système de coordonnées en radians, au moyen de \(2\pi\). Le changement
de coordonnées et son inverse sont
\begin{equation}
 \theta_A=\frac{\pi}{180}A_{\deg},\qquad
 \theta_C=\frac{\pi}{180}C^*_{\deg},\qquad
 A_{\deg}=\frac{180}{\pi}\theta_A,\qquad
 C^*_{\deg}=\frac{180}{\pi}\theta_C.
 \label{iv:angular:cambio-carta}
\end{equation}
Un parcours à nombre de tours conservé utilise les
coordonnées relevées. Le quotient circulaire exprime sa
classe modulo le tour correspondant.
\end{definition}

Le registre retient l’orientation, les tours et la mémoire du transport.
La phase finale et le relèvement sont des représentations différentes :
le second distingue des parcours qui partagent la première.
Chaque expression d’action déclare laquelle de ces coordonnées elle utilise.

\begin{proposition}[Covariance du lecteur d'action]
\label{iv:angular:accion}
Pour une section \(s\in\{\mathrm{pre},\mathrm{ret}\}\) et une
coordonnée angulaire relevée, on a
\begin{equation}
 \mathcal S_s(\theta)
 =\hbar_s\theta_{\mathrm{rad}}
 =h_s\frac{\theta_{\deg}}{360}.
 \label{iv:angular:accion-cartas}
\end{equation}
Un tour positif exprime \(h_s\).
\end{proposition}
\begin{proof}
Substituer
\(\theta_{\mathrm{rad}}=\pi\theta_{\deg}/180\) et
\(h_s=2\pi\hbar_s\) donne

\[
 \hbar_s\frac{\pi\theta_{\deg}}{180}
 =\frac{h_s\theta_{\deg}}{360}.
\]
En un tour, \(\theta_{\deg}=360\), d’où \(h_s\).

\end{proof}

La périodicité d’une représentation maintient sa condition
propre : si \(U(2\pi)=-I\), la restitution opératoire se produit
en \(4\pi\). La normalisation de l’action par tour de
\eqref{iv:angular:accion-cartas} et le retour de l’état
représenté conservent ainsi leurs domaines respectifs.

\subsection{Canaux positifs et réversion de la coordonnée impaire
}
\label{iv:angular:canales}

Les canaux sont exprimés, dans l’un ou l’autre des deux systèmes de coordonnées, au moyen de
\begin{equation}
 q_\pm=\exp[-(\theta_A\pm\theta_C)]
 =\exp\!\left[-\frac{\pi}{180}
                   (A_{\deg}\pm C^*_{\deg})\right].
 \label{iv:angular:q}
\end{equation}

\begin{proposition}[Récupération des coordonnées angulaires
]
\label{iv:angular:inversion}
Dans le système de coordonnées réel relevé, les canaux positifs déterminent
de manière unique
\begin{equation}
 A_{\deg}=-\frac{90}{\pi}\log(q_+q_-),\qquad
 C^*_{\deg}=\frac{90}{\pi}\log\frac{q_-}{q_+}.
 \label{iv:angular:recuperacion}
\end{equation}
L'échange \((q_+,q_-)\mapsto(q_-,q_+)\) conserve
\(A_{\deg}\) et inverse \(C^*_{\deg}\).
\end{proposition}
\begin{proof}
La somme et la différence des logarithmes de
\eqref{iv:angular:q} sont \(-\pi A_{\deg}/90\) et
\(-\pi C^*_{\deg}/90\), respectivement. Le logarithme réel
est univoque sur les canaux positifs. Leur produit est symétrique
et le quotient s’inverse lorsqu’on les échange.

\end{proof}

L’involution antipodale \(r\mapsto r+6\pmod {12}\) du
registre forme six paires. La distribution uniforme de la
coordonnée impaire sur ces paires attribue
\(C^*_{\deg}/6\) à chacune ; la somme récupère
\(C^*_{\deg}\). La conversion d’unités est appliquée une
seule fois à chaque argument. En particulier,
\(C^*_{\deg}/A_{\deg}=\theta_C/\theta_A\) lorsque
\(A_{\deg}\ne0\), tandis que l’état précédent conserve
les angles relevés et leur mémoire.
